\documentclass[11pt,a4paper]{article}
\usepackage[utf8]{inputenc}
\usepackage[T1]{fontenc}
\usepackage{lmodern}
\usepackage[portuguese]{babel}
\usepackage[a4paper,margin=2.5cm]{geometry}
\usepackage{xcolor}
\definecolor{TextEspresso}{HTML}{4A4238}
\definecolor{NudeTaupe}{HTML}{BFAFA6}
\definecolor{SoftBlush}{HTML}{D9C5B2}
\usepackage{titlesec}
\usepackage{enumitem}
\usepackage{listings}
\usepackage{hyperref}
\hypersetup{colorlinks=true, linkcolor=TextEspresso, urlcolor=TextEspresso}

\titleformat{\section}{\color{TextEspresso}\Large\bfseries}{\thesection}{0.5em}{}[{\color{NudeTaupe}\titlerule[0.8pt]}]
\titleformat{\subsection}{\color{TextEspresso}\large\bfseries}{\thesubsection}{0.5em}{}

\lstset{
  basicstyle=\ttfamily\small,
  breaklines=true,
  frame=single,
  rulecolor=\color{NudeTaupe},
  columns=fullflexible,
  showstringspaces=false
}

\newcommand{\guiaotitle}[2]{%
  \begin{center}
  {\Large\bfseries\color{TextEspresso} #1}\\[0.3cm]
  {\normalsize\color{NudeTaupe} #2}\\[0.2cm]
  {\color{NudeTaupe}\rule{4cm}{0.5pt}}
  \end{center}
  \vspace{0.5cm}
}

\begin{document}
\guiaotitle{Guião 13 -- WBEM com SFCB e Revisão do Semestre}{Infraestruturas de Computação na Nuvem, Aula 13}

\section{Objetivos}
\begin{itemize}[itemsep=0.2cm]
  \item Correr um CIMOM local (SFCB) num contentor Docker e consultá-lo com um cliente WBEM.
  \item Explorar classes CIM e a relação entre elas, contrastando com o modelo plano do SNMP.
  \item Consolidar, através de um roteiro de autoavaliação, os conceitos práticos de todo o semestre em preparação para o teste escrito.
\end{itemize}

\section{Pré-requisitos e Material}
\begin{itemize}[itemsep=0.2cm]
  \item Uma VM Linux no Proxmox (pode ser a \texttt{icn-lab} do Guião 01), com acesso \texttt{sudo}.
  \item Python 3 e \texttt{pip} instalados na VM.
  \item Acesso às pastas de todos os guiões anteriores (Aulas 01 a 12).
  \item Opcionalmente, credenciais de leitura da API do Proxmox, para a Parte 2.
\end{itemize}

\section{Parte 1 -- Explorar um CIMOM com WBEM}

\subsection{Introdução}
Vamos instalar o SFCB (Small Footprint CIM Broker), uma implementação leve de CIMOM, dentro da nossa VM, e consultá-lo a partir de um cliente Python baseado em \texttt{pywbem}. O objetivo é observar na prática a diferença entre o modelo orientado a objetos do CIM e o modelo plano de OIDs do SNMP visto na Aula 12 -- em particular, as \emph{associações}, que o SNMP não representa.

\subsection{Passo a Passo}
\begin{enumerate}[label=\arabic*.]

  \item \textbf{Instalar o CIMOM e os providers} na VM:
\begin{lstlisting}
$ sudo apt update
$ sudo apt install -y sfcb sblim-cmpi-base cim-schema wbemcli
$ sudo systemctl enable --now sfcb
$ sudo systemctl status sfcb
\end{lstlisting}
  O pacote \texttt{sblim-cmpi-base} fornece os \emph{providers} que traduzem entre o modelo CIM e o sistema real -- sem eles, o CIMOM não tem dados para expor.

  \item \textbf{Testar com o cliente de linha de comandos} \texttt{wbemcli}, antes de escrever código:
\begin{lstlisting}
$ wbemcli -nl ein 'https://root:<password>@localhost:5989/root/cimv2:CIM_OperatingSystem'
\end{lstlisting}
  A opção \texttt{-nl} desativa a verificação do certificado (autoassinado). Se a autenticação falhar, confirme a configuração de utilizadores em \texttt{/etc/sfcb/sfcb.cfg}; consoante a instalação, pode ser necessário usar as credenciais locais do sistema.

  \emph{Nota:} o SFCB e os seus providers têm empacotamento variável entre versões de distribuição. Se algum passo falhar, registe a mensagem de erro obtida -- a docente indicará a alternativa (outra imagem, ou dados de demonstração), e a análise do erro conta como parte do trabalho.

  \item \textbf{Instalar o cliente Python:}
\begin{lstlisting}
$ sudo apt install -y python3-venv
$ python3 -m venv ~/wbem-env
$ ~/wbem-env/bin/pip install pywbem
\end{lstlisting}

  \item \textbf{Criar o script} \texttt{explorar\_cim.py}, que liga ao CIMOM e enumera classes:
\begin{lstlisting}
import pywbem

conn = pywbem.WBEMConnection(
    'https://localhost:5989',
    ('root', 'PASSWORD'),
    default_namespace='root/cimv2',
    no_verification=True
)

print("=== Classes CIM disponiveis (root/cimv2) ===")
for cls in sorted(conn.EnumerateClassNames()):
    print(cls)
\end{lstlisting}
\begin{lstlisting}
$ ~/wbem-env/bin/python explorar_cim.py
\end{lstlisting}

  \item \textbf{Listar instâncias e propriedades} de uma classe de sistema:
\begin{lstlisting}
import pywbem

conn = pywbem.WBEMConnection(
    'https://localhost:5989',
    ('root', 'PASSWORD'),
    default_namespace='root/cimv2',
    no_verification=True
)

for inst in conn.EnumerateInstances('CIM_OperatingSystem'):
    print(f"--- Instancia: {inst.classname} ---")
    for prop, value in inst.items():
        print(f"  {prop}: {value}")
\end{lstlisting}
  Compare os nomes das propriedades com os OIDs numéricos usados na Aula 12: aqui são autodescritivos e tipados.

  \item \textbf{Navegar por associações -- o exercício central.} Esta é a operação que não tem equivalente no SNMP. A partir de uma instância, descubra os objetos que lhe estão associados:
\begin{lstlisting}
import pywbem

conn = pywbem.WBEMConnection(
    'https://localhost:5989',
    ('root', 'PASSWORD'),
    default_namespace='root/cimv2',
    no_verification=True
)

# Obter a instancia do sistema
sistemas = conn.EnumerateInstanceNames('CIM_ComputerSystem')
if not sistemas:
    raise SystemExit("Nenhuma instancia de CIM_ComputerSystem encontrada")

sistema = sistemas[0]
print(f"Sistema: {sistema}\n")

# Que associacoes partem deste objeto?
print("=== Associacoes (References) ===")
for ref in conn.ReferenceNames(sistema):
    print(f"  {ref.classname}")

# Que objetos estao associados a este?
print("\n=== Objetos associados (Associators) ===")
for assoc in conn.AssociatorNames(sistema):
    print(f"  {assoc.classname}")
\end{lstlisting}
  Registe as classes de associação encontradas. Para cada uma, procure perceber que relação do mundo real representa.

  \item \textbf{Invocar informação de uma classe} para ver os seus métodos, propriedades e chaves:
\begin{lstlisting}
import pywbem

conn = pywbem.WBEMConnection(
    'https://localhost:5989', ('root', 'PASSWORD'),
    default_namespace='root/cimv2', no_verification=True)

cls = conn.GetClass('CIM_OperatingSystem', IncludeQualifiers=True)
print("Propriedades-chave:")
for nome, prop in cls.properties.items():
    if 'Key' in prop.qualifiers:
        print(f"  {nome}")
print("\nMetodos definidos:")
for nome in cls.methods:
    print(f"  {nome}")
\end{lstlisting}
  Repare que a classe define \emph{métodos}, e não apenas valores -- outra diferença estrutural face ao SNMP.

  \item \textbf{Comparação direta com o SNMP.} Se ainda tiver o agente SNMP da Aula 12 acessível, obtenha o nome do sistema pelas duas vias e compare o esforço e a legibilidade:
\begin{lstlisting}
$ snmpget -v2c -c icnturma <IP> 1.3.6.1.2.1.1.5.0
\end{lstlisting}
  Registe as duas formas de chegar à mesma informação.

\end{enumerate}

\section{Parte 2 -- Gestão por API REST: o Proxmox (opcional)}
Esta parte liga a matéria da aula à tendência atual discutida na teoria: a substituição de SOAP/XML por APIs REST com JSON.

\begin{enumerate}[label=\arabic*.]
  \item \textbf{Consultar a API do Proxmox} a partir da VM, usando o token da Aula 07 (ou credenciais de leitura fornecidas pela docente):
\begin{lstlisting}
$ curl -k -H "Authorization: PVEAPIToken=<user>@pve!<token>=<segredo>" \
    https://<IP-do-servidor>:8006/api2/json/nodes
\end{lstlisting}

  \item \textbf{Consultar o estado de um node} e as suas métricas:
\begin{lstlisting}
$ curl -k -H "Authorization: PVEAPIToken=..." \
    https://<IP-do-servidor>:8006/api2/json/nodes/<node>/status
\end{lstlisting}

  \item \textbf{Comparar os três modelos de gestão} que usou ao longo do semestre para obter informação de um sistema:
\begin{lstlisting}
SNMP    -> snmpget sobre um OID numerico       (Aula 12)
WBEM    -> EnumerateInstances sobre uma classe (Parte 1)
REST    -> GET sobre um URL, resposta em JSON  (esta parte)
\end{lstlisting}
  Registe, para cada um, como se identifica o objeto a consultar e em que formato vem a resposta.
\end{enumerate}

\section{Parte 3 -- Roteiro de Revisão do Semestre}

\subsection{Introdução}
Esta segunda parte não introduz matéria nova: é um roteiro de autoavaliação para consolidar, de forma aplicada, os conceitos centrais de cada bloco do programa, como preparação direta para o teste escrito.

\subsection{Checklist por Bloco}
\begin{itemize}[itemsep=0.3cm]
  \item \textbf{C1 (Aulas 01-02)}: Rever a distinção entre os Tiers de datacenter e o respetivo SLA associado. Ser capaz de explicar, por palavras próprias, porque é que redundância de energia/rede aumenta o Tier.
  \item \textbf{C2 (Aulas 03-04)}: Comparar hipervisores tipo 1 vs tipo 2 e explicar porque é que um contentor partilha o kernel do host enquanto uma VM não. Rever o papel de namespaces e cgroups.
  \item \textbf{C3 (Aulas 05-07)}: Repetir mentalmente o ciclo de vida de um Pod no Kubernetes até ao self-healing. Explicar a separação control/data plane da SDN. Distinguir gestão de configuração (Ansible) de provisionamento declarativo (Terraform).
  \item \textbf{C4 (Aulas 08-09)}: Comparar scale-up vs scale-out e round-robin vs least connections. Explicar a diferença entre redundância ativa-ativa e ativa-passiva, e entre comunicação síncrona e assíncrona.
  \item \textbf{C5 (Aulas 10-11)}: Distinguir block, file e object storage com um exemplo de cada. Explicar RAID 1 vs RAID 5 e a diferença entre NAS e SAN. Enunciar o teorema CAP ao nível introdutório.
  \item \textbf{C6 (Aulas 12-13)}: Comparar a arquitetura manager/agente do SNMP com a arquitetura CIMOM/provider do WBEM. Identificar uma vantagem do WS-Management sobre o CIM-XML clássico.
\end{itemize}

\subsection{Exercício de Síntese}
Escolher \emph{um} cenário de infraestrutura (ex.\ ``uma aplicação web com picos de tráfego sazonais'') e, em poucas frases, descrever que tecnologias de cada bloco C1-C6 seriam relevantes para o desenhar e operar. Este exercício não é entregue, mas é fortemente recomendado como preparação para as questões de aplicação prática do teste escrito.

\section{Entregável / Questões de Verificação}
\begin{enumerate}[label=\alph*)]
  \setlength\itemsep{0.2cm}
  \item Enviar os scripts usados e a lista de propriedades obtidas para a instância escolhida.
  \item Submeta a lista de associações encontradas (\texttt{ReferenceNames} / \texttt{AssociatorNames}). Escolha uma e explique, por palavras próprias, que relação do mundo real representa.
  \item Porque é que uma navegação por associações como esta não tem equivalente direto no SNMP? O que teria de fazer para obter a mesma informação a partir de uma MIB?
  \item A classe que inspecionou define \emph{métodos}. Que diferença conceptual representa isto face às operações GET/SET do SNMP?
  \item Compare as três formas de obter o nome de um sistema que usou no semestre (SNMP, WBEM e, se fez a Parte 2, REST). Como se identifica o objeto em cada caso e em que formato vem a resposta?
  \item Que vantagem prática (não apenas teórica) sentiu ao consultar um CIMOM com \texttt{pywbem} em comparação com o \texttt{snmpwalk} da Aula 12? E que desvantagem?
  \item Se tivesse de escolher hoje uma tecnologia para gerir uma frota de servidores heterogéneos, qual escolheria e porquê? Considere SNMP, WBEM/WS-Management e APIs REST.
\end{enumerate}

\end{document}
